\documentclass[../../../../SoftwareRequirementsSpecification.tex]{subfiles}
\begin{document}
% Richiede le macro definite in src/macros/useCaseMacros.tex
% (incluse nel preambolo di SoftwareRequirementsSpecification.tex)

\ucstart
\begin{xltabular}{\textwidth}{|l|X|}
  \hline
  \ucid{PT-12} \\ \hline
  \ucname{Registrazione Palestra} \\ \hline
  \ucactor{PT} \\ \hline
  \ucdescription{Il PT registra nel sistema una palestra non ancora
    presente, per potersi in seguito associare ad essa (caso d'uso
    PT-13).} \\ \hline
  \ucprecond{Il PT deve essere autenticato.} \\ \hline
  \ucpostcond{Il sistema crea la palestra. Nessun PT o atleta è ancora
    associato ad essa.} \\ \hline
  \ucmainflow{
    \item Il PT apre la pagina di ricerca palestra (vedi Nota 1).
    \item Il PT preme il pulsante "Registra nuova Palestra".
    \item Il sistema mostra il form di inserimento dati della
      palestra (nome, indirizzo e telefono).
    \item Il PT inserisce i dati.
    \item Il PT preme il pulsante "Registra Palestra".
    \item Il sistema verifica che i dati siano validi.
    \item Il sistema crea la palestra.
    \item Il sistema mostra una pagina di conferma.
  } \\ \hline
  \ucaltflow{
    \textbf{6a. Dati non validi} \newline
    - Il sistema mostra un messaggio di errore ("Controllare i dati
    inseriti") e indica i campi non validi. \newline
    - Il flusso riprende dal passo 4.
  } \\ \hline
  \ucnote{
    \textbf{Nota 1:} la pagina di ricerca palestra è la stessa usata
    al passo 1 del caso d'uso PT-13: il PT vi cerca la palestra a cui
    associarsi, e solo se non la trova può registrarla da qui. Questo
    evita di creare più volte la stessa palestra nel database.
    \newline
    Il sistema non verifica l'esistenza reale della palestra: si
    limita a registrare i dati che il PT dichiara.
  } \\ \hline
\end{xltabular}
\end{document}
